% fhe_full_transcript.tex
% Full verbatim transcript of Craig Gentry's EUROCRYPT 2021 talk on FHE
% Compile: xelatex fhe_full_transcript.tex

\documentclass[11pt,a4paper]{article}

\usepackage{fontspec}
\usepackage{xeCJK}
\setCJKmainfont[BoldFont=PingFang SC Semibold]{PingFang SC}
\setCJKsansfont[BoldFont=PingFang SC Semibold]{PingFang SC}
\setCJKmonofont{PingFang SC}
\setmainfont{Palatino}
\setsansfont{Helvetica Neue}
\setmonofont{Menlo}

\usepackage[top=2.5cm, bottom=2.5cm, left=2.5cm, right=2.5cm]{geometry}
\usepackage{amsmath,amssymb}
\usepackage{graphicx}
\graphicspath{{images/}}
\usepackage{enumitem}
\usepackage{xcolor}
\usepackage{hyperref}
\usepackage{caption}
\hypersetup{
  colorlinks=true,
  linkcolor=blue!70!black,
  urlcolor=blue!70!black,
  pdfauthor={Whisper Transcription},
  pdftitle={A Decade (or So) of Fully Homomorphic Encryption - Full Transcript},
}

\usepackage{parskip}
\setlength{\parindent}{0pt}
\setlength{\parskip}{6pt plus 2pt minus 1pt}

\usepackage{mdframed}
\newmdenv[
  topline=false, bottomline=false, rightline=false,
  linewidth=3pt, linecolor=blue!40,
  innerleftmargin=12pt, innerrightmargin=12pt,
  innertopmargin=8pt, innerbottommargin=8pt,
  backgroundcolor=blue!3, skipabove=10pt, skipbelow=10pt,
]{myquote}

\usepackage{titlesec}
\titleformat{\section}{\Large\bfseries\sffamily}{}{0em}{}[\vspace{-0.5em}\rule{\textwidth}{0.5pt}]
\titleformat{\subsection}{\large\bfseries\sffamily}{}{0em}{}
\setcounter{secnumdepth}{0}

\setlength{\emergencystretch}{3em}
\tolerance=1000

\newcommand{\keyslide}[2]{%
  \begin{center}
    \includegraphics[width=0.88\textwidth]{#1}
    \captionof{figure}{#2}
  \end{center}
  \vspace{6pt}
}

\newcommand{\timestamp}[1]{%
  \hfill{\scriptsize\color{gray}[#1]}%
}

\begin{document}

\begin{titlepage}
  \centering
  \vspace*{2.5cm}
  {\Huge\bfseries\sffamily A Decade (or So) of\\Fully Homomorphic Encryption\par}
  \vspace{0.8cm}
  {\LARGE 全同态加密十（几）年的发展历程\par}
  \vspace{1.2cm}
  {\Large\sffamily Craig Gentry\par}
  {\normalsize Algorand Foundation\par}
  \vspace{0.5cm}
  {\normalsize EUROCRYPT 2021 Invited Talk\par}
  \vspace{0.3cm}
  {\normalsize 翻译/字幕：刘巍然（学酥）\par}
  \vspace{1.5cm}
  {\large\bfseries 完整逐字稿（Full Verbatim Transcript）\par}
  \vspace{0.3cm}
  {\small Transcribed by Whisper large-v3-turbo\par}
  \vspace{0.5cm}
  {\small 原始视频：\href{https://www.bilibili.com/video/BV1rY411V7Ko/}{Bilibili BV1rY411V7Ko}\par}
  \vfill
  {\large \today\par}
\end{titlepage}

\newpage
\tableofcontents
\newpage


\section{Introduction \& FHE Origins (00:00 -- 10:00)}

\keyslide{slide_0001.jpg}{Slide at 00:00}

And we are very happy and very honoured to have Craig Gentry as invited speaker today. So Craig is a research fellow at the Algorand Foundation. He received his PhD from Stanford in 2009. And as you all know, he wrote an amazing dissertation in which he constructed the first FHE scheme. And so his dissertation, of course, was awarded. \timestamp{00:00}

\keyslide{slide_0002.jpg}{Slide at 00:39}

He won many awards, including the ACM Doctoral Dissertation Award and the Grasshopper Award. And so Craig is also the author of many major results on HFE. For instance, he's the G letter in BGV. He's the G letter in GSW. And he has also important contribution on zero-knowledge proofs, on aggregate signatures, and on obfuscation. \timestamp{00:36}

\keyslide{slide_0003.jpg}{Slide at 01:19}

So today, his talk is entitled A Decade or So of Fully Homomorphic Encryption. So please join me to welcome Craig Gentry. Thank you so much for that great introduction and for inviting me. So I'll be talking about the past decade or so of Fully Homomorphic Encryption. But of course, the story starts much earlier than that. \timestamp{01:13}

\keyslide{slide_0004.jpg}{Slide at 01:58}

That's the advent of public key cryptography with these guys, Ravest, Edelman, and Tertuzos, who noticed that the RSA scheme had a very peculiar property, that basically you could take the ciphertexts and multiply them. And that would induce multiplication of the underlying messages. And you could do that, of course, without the secret key. So in effect, you could compute on the encrypted data while it remains encrypted. And that allows you to separate the processing of data and the access of data. \timestamp{01:49}

\keyslide{slide_0005.jpg}{Slide at 02:38}

And in fact, they had, back in 1978, they had this idea of basically what we would call today encrypted cloud computing. So there you have a user, Alice, who wants to use cloud computing, but wants to keep her data secret. And so, of course, she encrypts her data. And later, she has a query about her data. She wants to retrieve certain files from the cloud that contain some keyword or something like that. \timestamp{02:26}

And what she wants is that the cloud, despite the data being encrypted, should be responsive to the query. It should be able to perform the query. So what you need is for the encryption scheme to have sort of a magical algorithm by which it can take that function that Alice sent and her encrypted data and evaluate it and produce a ciphertext that encrypts the function applied to the data inside that encryption. So that's what a homomorphic encryption scheme does. \timestamp{03:02}

\keyslide{slide_0006.jpg}{Slide at 03:18}

It runs that evaluation function. And at that point, the cloud can send back the encryption of the response that Alice wants, and she can decrypt it. What we want from homomorphic encryption is that at no point does the cloud see any information about Alice's data at any intermediate stage of the computation. Nothing at all should leak about the data. I mean, no partial information or anything like that. \timestamp{03:39}

\keyslide{slide_0007.jpg}{Slide at 03:57}

The only thing that should leak about the data is the length of the data, because, of course, the ciphertext length has to be an upper bound on the length of the data, right? So that much will leak, but aside from that, nothing at all should leak in this scheme. And her data should be secure, not only against the cloud, but, of course, against any adversaries that breach the cloud or against insider attacks, and that sort of thing. So, okay, so in 2009, I proposed the first plausibly secure fully homomorphic encryption scheme. Sorry about the picture. \timestamp{04:16}

\keyslide{slide_0008.jpg}{Slide at 04:37}

This is the only good one I could find. So basically, the idea of that scheme was sort of a self-referential thing that we call bootstrapping, and that sort of relates to the question of what if a homomorphic encryption scheme could decrypt itself with an encrypted secret key? So it takes the encrypted data of the secret key and sort of decrypts itself and sort of this loop, Ouroboros loop that I had on the first slide. And it turns out that that idea is very powerful. It's allowed us to get fully homomorphic encryption schemes and much better ones than I proposed in 2009, certainly. \timestamp{04:52}

\keyslide{slide_0009.jpg}{Slide at 05:16}

So this scheme was very slow. You might wonder how slow it could be. So slow enough that it prompted Butler Lamson to say, I don't think we'll see anyone using Gentry's solution in our lifetimes. So that kind of hurt. But in retrospect, he's completely right, because the current solutions that we have now are so much better by a factor of, like, literally millions. \timestamp{05:29}

\keyslide{slide_0010.jpg}{Slide at 05:56}

And lots of people have contributed. We have now four generations of fully homomorphic encryption schemes. We have standardization happening. So it's been an amazing journey. So here's the organization of the rest of the talk. \timestamp{06:05}

\keyslide{slide_0011.jpg}{Slide at 06:36}

So first, I want to talk. You have the arrow of time there. And so first, I want to talk about the past and past 10 years, where it's gone from a proof of concept to, I think it's fair to say, a usable tool. I'll kind of barrage you with, you know, some performance numbers and, you know, all sorts of applications. Even though I'm sort of a theoretical sort of person, I'm not the sort of person that writes, that prototypes, you know, FHE evaluation of neural nets. \timestamp{06:42}

\keyslide{slide_0012.jpg}{Slide at 07:15}

So maybe I'm not the best proponent of this, but I feel like it's kind of my responsibility to update people's belief on the performance of FHE, given that I was maybe a little too honest about the performance of FHE back in 2009, and people's beliefs have maybe sort of fixed at that point. So I think we kind of need to update our beliefs and look at the performance as it is now. So for the present or sort of timeless part, I want to talk, I'm going to describe how the current FHE schemes work, but I want to do it in a different way, partly because I'm sort of bored of the typical, you know, lattice, LWE, blah, blah, blah, presentation of FHE. That sort of, you know, that sort of, you know, that sort of, you know, that sort of, you know, that sort of, you know, are soulmates, you know, in some sort of perfect union or something. So I kind of want to divorce it from Lattice's a little bit and sort of take a mathematician's point of view and start from the homomorphism and then kind of describe how to sort of harden or mask the homomorphism to create a semantically secure encryption scheme. \timestamp{07:19}

\keyslide{slide_0013.jpg}{Slide at 07:55}

And then in the final part, I want to talk about the future, you know, is there any hope of sort of escaping the blueprint for FHE that we have now, which uses bootstrapping and lattices and noisy cipher texts? Nothing would make me happier than to sort of get rid of all that, you know, all that baggage and create, you know, fundamentally new schemes. That would be very exciting. Okay. So now for the first part where I talk about how it's become a usable tool over the past 10 years. \timestamp{07:55}

\keyslide{slide_0014.jpg}{Slide at 08:35}

So, as I said, there have been four generations of FHE. The first generation started with my initial paper and then we had a scheme based on integers and kind of a simplification. And then we actually implemented the first scheme, including the bootstrapping, which is the most complicated step back in 2011. But the performance was truly atrocious. And then the second generation started with this paper of Pekersky and Bakertanathan, where they showed how to base FHE on the learning with errors problem, LWE, which is a well-established cryptographic assumption. \timestamp{08:32}

\keyslide{slide_0015.jpg}{Slide at 09:14}

So that was very nice. But they also, you know, so in addition to that security benefit, there was also an efficiency benefit that they had in there. They had a technique called modulus reduction and in a follow-up paper, which is usually called BGV, we took modulus reduction to the max and created a much more efficient scheme. And so, you know, so that's a good performance improvement, which was a good performance improvement, which is, you know, a lot of the parameters sort of grew up the depth of the circuit that you're evaluating, as opposed to the degree, it's an exponential difference between the depth and the degree. So it's a, it was a good performance improvement. \timestamp{09:08}

\keyslide{slide_0016.jpg}{Slide at 09:54}

Another part of the second generation is that we had SIMD operations. Basically, you can pack like a thousand different. \timestamp{09:45}


\section{Four Generations of FHE (10:00 -- 20:00)}

plaintext inside each ciphertext in little 1,000 different slots and operate on these slots in parallel. And on top of that, you can also, there's a way that you can use automorphisms to sort of permute the different slots so that they can interact with each other. And so back in 2012, we had a scheme which actually had polylog overhead, polylog in the security parameter, which was really amazing because you would expect that the overhead of any scheme would be at least \timestamp{10:00}

\keyslide{slide_0017.jpg}{Slide at 10:33}

the security parameter, right? But it's only polylog in the security parameter. So we thought we were done at that point in terms of performance because, I mean, how can you be better than polylog? Didn't turn out that way. So the performance did improve incredibly over those initial years. So you can see here, this doesn't include bootstrapping, so that's a little qualification here. These are kind of moderate depth circuits, but we could perform on the order of a thousand, \timestamp{10:38}

\keyslide{slide_0018.jpg}{Slide at 11:13}

a few thousand gates per second. You can see we're beating Moore's law by a huge margin. This is logarithmic scale. And so that was good, but this doesn't include bootstrapping. So if you want to see where we are in 2020 for like generation two schemes, this is what it looks like. This is bootstrapping in Helib. And if you look at the sort of the bottom highlighted line there, you can see that the amortized time per plain text bit for bootstrapping is 0.9 milliseconds. So that's pretty good. I mean, the actual time to bootstrap a ciphertext itself is 11 seconds, which you can see on the line right above it, but each ciphertext contains a lot of different plain text bits. \timestamp{11:16}

\keyslide{slide_0019.jpg}{Slide at 11:53}

And on top of that, you don't have to bootstrap at every level of the arithmetic circuit that you're evaluating. It says right there that there are 17 usable levels. So you, every 17 levels will bootstrap. So amortized, you get that it's 0.9 milliseconds per plain text bit. So, uh, so that's much better. Um, you'll see what the third and fourth generation schemes that we can do even, even better than that. Okay. So the third generation, um, I guess kind of started with a paper that I had with Mitzahai and Redwaters, uh, which was kind of a simplification. And also it improved the efficiency a little bit of, uh, of prior schemes. \timestamp{11:55}

\keyslide{slide_0020.jpg}{Slide at 12:32}

But what was really cool, uh, was a follow-up paper by Rukursky and by Gutanathan, where they showed that they, you could take that scheme and it had a very nice property that you could sort of tweak it. And, um, uh, you could get very good parameters. So the noise wouldn't grow very much. So the basically you could get a security level just as good as regular lattice-based PKE. Okay. Normally FHE has very like huge parameters, which entail like a big loss in security, but here they showed that you could actually make it as secure as PKE. And on top of that, they showed that, uh, there was a very, well, related to that, there was, um, you could tweak the scheme to get a very efficient bootstrapping technique, which, uh, made it more efficient. \timestamp{12:33}

\keyslide{slide_0021.jpg}{Slide at 13:12}

And so in a series of follow-up papers, the bootstrapping became more and more efficient. So, right, you know, the last paper, it's, uh, less than a tenth of a second, and that was back in 2016. So I, I think, uh, the performance has probably improved since then. Um, so for the third generation, the, the bootstrapping is faster. Um, but one sort of downside of these schemes is that, uh, they don't natively support ciphertext packing. \timestamp{13:12}

\keyslide{slide_0022.jpg}{Slide at 13:51}

So in the second generation, you could stick tons of plain text inside each ciphertext. That's not sort of natively supported in this, in this, uh, third generation. The fourth generation, uh, was the CKKS, uh, scheme, which, uh, operates on floating point numbers, which is very useful for real world, uh, applications like neural networks. So that, uh, it really, uh, helped efficiency immensely. Um, so here are libraries for FHE. \timestamp{13:50}

\keyslide{slide_0023.jpg}{Slide at 14:31}

You can see, well, first of all, there are lots of libraries. And, and secondly, most of them, uh, implements, uh, CKKS because it's, uh, so useful, but the other generations are, are also useful too. And in fact, there's something called, uh, chimeric FHE, where basically the idea is that, uh, it's very useful to actually switch in between different FHE schemes. So, um, so FHE, but, you know, using the same process as bootstrapping, where basically you apply a decryption function homomorphically, you can bootstrap out of one FHE scheme into another. So you can sort of translate between different FHE schemes and, uh, you know, depending on the context, maybe you want a third generation scheme that operates very well on bits or a fourth generation scheme that operates very well on floating point numbers. \timestamp{14:29}

\keyslide{slide_0024.jpg}{Slide at 15:11}

So you can switch in between them. So, uh, this is a slide, uh, courtesy of, uh, uh, uh, Jung-hee Chion, um, who's done a lot of the, the, the recent performance improvements. Um, he shows how the, the bootstrapping time has, uh, improved over the years. Um, right now they have, um, on a CPU based implementation, they have, uh, 19, uh, microseconds per bit, uh, amortized bootstrapping time. So there you see that the ciphertexts are kind of big. \timestamp{15:07}

\keyslide{slide_0025.jpg}{Slide at 15:50}

I mean, they're 262 kilobytes, but they hold lots of, uh, different plain texts and you only have to bootstrap every once and however many levels. So you get, uh, it only takes 35 seconds to bootstrap that, uh, really big ciphertext. So you get, uh, 19 microseconds amortized. Um, and then they also did a GPU based implementation and there it's a sub microsecond per bit amortized bootstrapping time, which is just, uh, just amazing. Okay. \timestamp{15:46}

\keyslide{slide_0026.jpg}{Slide at 16:30}

Okay. Um, so now I'll talk a little bit about some of the applications again, I want to, you know, sort of clarify that I'm, you know, I'm kind of a theoretical, theoretical guy. So, um, you know, uh, I don't know all of the details of these implementations. Um, so I sort of kind of sifted through them to see kind of the most important and most, uh, impressive ones. Um, so if that qualification, uh, uh, I'll describe a few, uh, so one, um, big application area for FHE has been, uh, genomic analysis and that's been in part because there's this consortium called IDASH, uh, that every year runs some, uh, challenges, um, you know, relating to MPC and, and FHE trying to, uh, run privacy preserving algorithms over genomic data. \timestamp{16:24}

\keyslide{slide_0027.jpg}{Slide at 17:10}

Um, and I guess cryptographers love a challenge. So every year there are, uh, several teams that, uh, that enter this competition and produce pretty impressive results. Um, so here, uh, I guess this is, uh, 2019 they did, uh, genome wide association studies and, uh, excuse me, genotype, uh, imputation. Um, and they, the, the types of algorithms they used here are like linear regression, logistic regression, and neural networks. And for 80,000 SNPs, they were able to do, uh, an evaluation time of, uh, 27, uh, 25 seconds. \timestamp{17:03}

\keyslide{slide_0028.jpg}{Slide at 17:49}

And there was also some, uh, some training aspect, uh, uh, to this as well. Uh, another big application area as neural networks, uh, probably the first paper to really attempt, uh, sort of a, uh, a moderate depth, uh, uh, neural network with the, uh, proper nonlinear activation functions was this, uh, Kryptonets paper, uh, from 2016. Uh, since then, there've been, uh, tons of, uh, papers on neural net. \timestamp{17:41}

\keyslide{slide_0029.jpg}{Slide at 18:29}

And, um, I just picked, uh, an example from this year that I really liked where they, um, uh, they ran, uh, speech recognition on, um, um, to, um, distinguish among a hundred different people, and they managed to accomplish that in less than one second. Without, without dumbing down the speech recognition algorithm at all, just running it as is, they're able to distinguish among a hundred people in less of a second. Uh, another application is a private information retrieval. Uh, there, the idea is that you would like to retrieve one item from a database without telling the server which item you're retrieving. \timestamp{18:20}

\keyslide{slide_0030.jpg}{Slide at 19:08}

Um, so one, um, sort of, uh, necessary drawback of private information retrieval is that the server has to touch every item in the database. Because if it didn't touch a particular item, you would know somehow that, that you weren't interested in that particular item. So by necessity, the server has to touch every item in the, in the database when it responds to a, uh, a peer query. Uh, so server overhead is a big concern because it's touching the whole database. And if the, you know, cryptographic overhead on top of that is huge, then that would be a bad thing. \timestamp{18:58}

\keyslide{slide_0031.jpg}{Slide at 19:48}

So Sion and, uh, Carbonar had this paper back in 2007, where they said, there's really no hope for single server private information retrievals to ever be, uh, more efficient than just the server just sending the entire database to the, to the user. Uh, and their analysis was based on number of theoretic schemes like Pallier, um, and for those, \timestamp{19:36}


\section{Performance, Usability \& Standardization (20:00 -- 30:00)}

maybe that's true, but shortly after that, sort of a re-analysis was done based on a lattice-based peer scheme, which was orders of magnitude faster, saying that this analysis no longer applies. For the server, it's going to be more efficient to actually respond to the peer query than to just communicate the entire database. So peer schemes have become faster and faster. I'll mention one paper that I had with Shai Levy, which is super fast, and the ratio between the plain \timestamp{20:00}

\keyslide{slide_0032.jpg}{Slide at 20:28}

text and the ciphertext is four over nine. So you only have like a factor of two ciphertext expansion. So the communication bandwidth is practically optimal, assuming a certain lower bound on the amount of information that the courier wants. And we didn't do an implementation ourselves, but there's a preliminary implementation by Samir Manon and David Wu. And what they found was that they could do a peer query on a million files of 30 kilobytes each in 86 seconds. \timestamp{20:36}

\keyslide{slide_0033.jpg}{Slide at 21:07}

You can compare that to unaccelerated AES encryption of the whole database, and that takes about the same amount of time, 85 seconds. So the amount of time for the server to respond to a peer query is basically the same as a whole database encryption using AES, which I think is really amazing. Okay. So that, you know, perhaps you think that comparison was a little unfavorable in the sense that this uses unaccelerated AES, but you can also hardware accelerate FHE, and there are big \timestamp{21:12}

\keyslide{slide_0034.jpg}{Slide at 21:47}

efforts in that regard. There's a DARPA program with 70 million funding over four years, whose objective is basically to produce good hardware accelerators for FHE. So Intel and Microsoft are in there, and then there are a few other groups. But of course, as I mentioned, the hardware acceleration is already happening without this program. So this paper out of Korea that I mentioned earlier, Junghee's group, here in the paper, they have 0.423 microseconds per bit amortized. So that's actually, I mean, \timestamp{21:48}

\keyslide{slide_0035.jpg}{Slide at 22:26}

and then he just communicated to me recently that it's down to 0.29. So I mean, so even in this, this paper was actually posted this year. So just within the past few months, he's improved it by another factor of by 25\% or so. So it's, it's almost as if you can see the performance improving in real time. All right. So one important factor for FHE is usability. There are all sorts of different toolkits that help you translate your program on unencrypted data into one-on encrypted data. \timestamp{22:24}

\keyslide{slide_0036.jpg}{Slide at 23:06}

So Google recently made the news by open sourcing a transpiler that, you know, translates high-level code into code that operates on encrypted data. The idea being that right now, you kind of need to be a crypto a crypto expert to do this, these transformations. But I haven't looked closely at this transpiler, but the idea is that basically not much crypto expertise would be needed. In terms of standardization, first of all, NIST has this standardization of post-quantum crypto going on. \timestamp{23:00}

\keyslide{slide_0037.jpg}{Slide at 23:46}

They started in 2015 with 82 submissions. And post-quantum crypto is crypto that we hope, think, will be secure after large-scale quantum computers are built, you know, when that happens. And out of the seven finalists, you'll see that five of them are lattice-based. So lattice-based crypto certainly is on its way to being standardized. It's here to stay. And FHE is going to be a part of that. And I'll mention again that the standardization workshop for home-work encryption, they have a \timestamp{23:36}

\keyslide{slide_0038.jpg}{Slide at 24:25}

draft standard for various home-work encryption schemes of different generations. They have ongoing white papers on security levels, APIs, and applications. So you can go there and look at that. Okay. So now I'm going to go on to the second part of my presentation, where I'm going to be talking about how current home-work encryption schemes are built. But I'm going to do it from the standpoint of kind of a mathematician who might \timestamp{24:13}

think first about a homomorphism and think second about how can I sort of tweak this to make it secure. Okay. So first of all, I haven't provided any definitions yet. So I should probably define some things. So here's a public key encryption. You have the usual key generation encryption and decryption algorithms. So key generation and encryption are probabilistic algorithms. So I should probably put a little R inside as one of the inputs to the encryption algorithm for randomness. But I didn't. But just keep in mind that \timestamp{24:49}

\keyslide{slide_0039.jpg}{Slide at 25:05}

it's probabilistic. Correctness is just what you would think that, you know, if you encrypt the message M and then decrypt that ciphertext, you should get back M for all key pairs, messages, and encryption randomness. And semantic security is kind of the right notion of security for an encryption scheme. And it was invented by Goldwasser Macaulay in 1982. And what it says for our purposes is basically that the encryptions of zero and encryptions of one should be indistinguishable. It should be hard to distinguish \timestamp{25:25}

\keyslide{slide_0040.jpg}{Slide at 25:45}

for computationally bounded adversaries. Okay. So any secure public encryption scheme must be probabilistic. There must be many ciphertext per message because if they were just, you know, one ciphertext per message, you know, if they're one encryption of zero and one encryption of one, well, it'd be pretty easy for you to break the scheme then, right? Because especially since it's public key, you would just encrypt zero and compare it to the \timestamp{26:01}

\keyslide{slide_0041.jpg}{Slide at 26:24}

ciphertext you have in your hand. And then you would be able to tell what the ciphertext in your hand encrypts. So any PKE scheme must be probabilistic. So for homewarf encryption, you have a fourth algorithm in addition to the first three, which we're called V for evaluate. And basically what it should do is, well, it's kind of like that initial picture I had with encrypted cloud computing. Basically, if you have a bunch of ciphertexts, c1 through ct, which encrypt m1 through mt, if you plug them into \timestamp{26:37}

\keyslide{slide_0042.jpg}{Slide at 27:04}

the evaluate function together with some function f, what the output should be is a ciphertext that encrypts f applied to m1 through mt. Okay. And security is defined the same way as basic PKE, in other words, semantic security. And so again, there should be many ciphertexts per message. So it might be a little easier to see what homomorphism means in terms of this, what's called a commutative diagram. So you start off in the upper left corner with a bunch of ciphertexts. C stands for \timestamp{27:13}

\keyslide{slide_0043.jpg}{Slide at 27:43}

ciphertexts. And you could just take the downward arrow and just decrypt those two ciphertexts and get team messages. And then at that point, apply f to the team messages to get some other message. Or you could sort of take the northeast passage and go the rightward arrow and apply the evaluate function to the t ciphertexts using the function f inside to get a ciphertext that, when decrypted, will give you that message. And the commutative diagram, the main property is that it doesn't \timestamp{27:49}

\keyslide{slide_0044.jpg}{Slide at 28:23}

matter which way you go, you know, whether you take the northeast passage or the southwest passage, either way you should get the same thing. So it's commutative because the order in which you apply the decryption function and the function f don't matter. It doesn't matter. You can do it either way. You can take the northeast passage and apply the function f and then decrypt, or you can just immediately decrypt and then apply the function f. You should get the same value. \timestamp{28:26}

\keyslide{slide_0045.jpg}{Slide at 29:03}

Okay, so every homomorphic encryption scheme has a family of functions associated to it, which is the family of functions that it actually works for, that it's correct for. One family of functions could be the family of arithmetic circuits, basically layered graphs of addition and multiplication gates. So you can express boolean circuits like this, and boolean circuits are very powerful. Basically everything can be expressed as a boolean circuit. \timestamp{29:02}

\keyslide{slide_0046.jpg}{Slide at 29:42}

So if your homomorphic encryption scheme can do everything, can do all functions, and we call it a fully homomorphic encryption scheme. So that's what we're aiming for. Okay, so now there are kind of two ways. They're all \timestamp{29:38}


\section{Two Approaches to Homomorphic Encryption (30:00 -- 40:00)}

two ways of viewing things, right? So for the purposes of this presentation, I'm saying there are two ways to develop homework and encryption schemes. The first way is the righteous way, the cryptographic way, exemplified by Goldwasser and Macaulay, where you start with an established cryptographic assumption, you build public key encryption on the assumption, and then at that point, maybe you notice that the scheme is already homomorphic, or you try to try to make it \timestamp{30:00}

\keyslide{slide_0047.jpg}{Slide at 30:22}

homomorphic by introducing some sort of addition and multiplication operations. The pro of this approach is that you're less likely to get an insecure scheme, because you started with a well-established cryptographic assumption. The con is that you're not likely to get any homomorphic encryption scheme at all, unless you were just kind of lucky, and your public key encryption scheme that you created was already homomorphic. The alternative way that I'd like to \timestamp{30:42}

\keyslide{slide_0048.jpg}{Slide at 31:01}

suggest here, exemplified by Galois here, is the mathematical way, where you start with a set of homomorphisms. And you build, you kind of interpret decryption as belonging to that family of homomorphisms, and you build the homomorphic encryption scheme off of that. And then at that point, you start thinking about security, like how can I sort of harden or patch this, or mask this homomorphism, so that it becomes secure, so that you can get semantically secure encryption from it. So how do I make my ciphertext look \timestamp{31:25}

\keyslide{slide_0049.jpg}{Slide at 31:41}

\keyslide{slide_0050.jpg}{Slide at 32:21}

indistinguishable from random? So the pro of this approach is that it sort of allows you to look at different mathematical structures and consider whether they would be useful for cryptography. The con is, well, your scheme is likely to be broken. And then you're, I don't know, then you're one of those people that proposes lots of broken schemes, and you have to suffer eternal shame. \timestamp{32:08}

\keyslide{slide_0051.jpg}{Slide at 33:00}

So I don't want to propose that the mathematical way is the right way. I think both of these ways are useful, even internally within the same person, I think you should at times consider both approaches. But I certainly don't want to die in a duel with cryptographers, similar to Galois here. Okay, so what is a homomorphism? Well, basically, it's a structure preserving map between two algebraic structures that preserves the operations. \timestamp{32:51}

\keyslide{slide_0052.jpg}{Slide at 33:40}

And it's useful to draw these things with commutative diagrams. So here's an example of a homomorphism, just taking x to e to the x. So in this case, if you have x and y, you start off with x and y in the upper left corner, you could follow the right arrow and add them together, you get x plus y. And then you apply delta, you go downward, and you get e to the x plus y. Or starting with x plus y in the upper left, you could go down immediately, you get e x and e y, and you multiply them together, and you get e to the x plus y. Okay, so that's just an example of homomorphism. So notice that the operations in the two structures don't have to be the same. So for the first structure, it's a plus, and for the second, it's times. \timestamp{33:34}

\keyslide{slide_0053.jpg}{Slide at 34:20}

Okay, so where things become interesting is when you combine homomorphisms with complexity theory. So let's just look at a generic commutative diagram. And let's imagine that the downward arrows are the homomorphisms, and the rightward arrows are like operations. Then, now let's suppose that the rightward arrows are easy to compute, they're tractable. Whereas the downward arrows, the homomorphisms are hard to compute. \timestamp{34:17}

\keyslide{slide_0054.jpg}{Slide at 34:59}

Unless you have some trapdoor, secret key, or something like that. Well, that's basically what homomorphic encryption is. It's just a very generic combination of homomorphisms with complexity theory. Basically, you can view the top level of this commutative diagram as the ciphertext level, where you're just, you know, you have ciphertexts in A, B, and C, and you can perform these operations at that level. \timestamp{35:00}

\keyslide{slide_0055.jpg}{Slide at 35:39}

But it's hard to go take a downward arrow down to the plaintext level. So, as I said, homomorphic encryption is basically homomorphisms plus complexity theory. So, all right. So this commutative diagram sort of explained what's happening with the homomorphisms and the decryption. The decryption is homomorphism. \timestamp{35:42}

\keyslide{slide_0056.jpg}{Slide at 36:18}

But what about keygen and encrypt? So, Rothblum has this nice result where you can take a semantically secure private key homomorphic encryption scheme that can do homomorphically edition modulo 2, and just generically convert it to, you know, subject to some qualifications, convert it to a semantically secure public key encryption, homomorphic encryption scheme that can do the same operation homomorphically. \timestamp{36:25}

\keyslide{slide_0057.jpg}{Slide at 36:58}

And the technique is basically this. So, as your public key, you just put a lot of encryptions of 0 and encryptions of 1. And then in the encryption process, you just apply homomorphic addition among a random subset of the ciphertext in your public key to produce a new ciphertext, which encrypts your message. A little more specifically, you have k random encryptions of both 0 and 1 each, call them c's and d's. \timestamp{37:08}

\keyslide{slide_0058.jpg}{Slide at 37:38}

And then in encryption, you pick a random subset of the numbers 1 through k. Such that your message is equal to the cardinality of S modulo 2. And then you just homomorphically add the ciphertext corresponding to the ciphertext corresponding to S that encrypt one and use the encryptions of 0 for the complement of S. So, you can see that if you add these ciphertexts, the parity of what is encrypted is equal to the number of d's, which is the cardinality of S. \timestamp{37:51}

\keyslide{slide_0059.jpg}{Slide at 38:17}

OK. So, this is just a generic way that you can produce public key homomorphic encryption from private key homomorphic encryption. And it's nice because it also simplifies proving security of your scheme because this encryption procedure he proves is secure and all you have to worry about at this point is proving the indistinguishability of the ciphertext that you put as part of the public key. So, it simplifies provable security. So, in that sense, the homomorphic way of producing public key encryption schemes is kind of compatible in spirit with the cryptographic way. OK. So, now I'm going to give you kind of a simple example of how Goldwasser Macaulay, the encryption scheme, might have been built by someone who was looking at it the mathematical way. So, I don't know exactly how Goldwasser Macaulay did it, but you can imagine as a sort of a counterfactual example, this might have happened that you start with a homomorphism, namely the Legendre symbol. \timestamp{38:34}

\keyslide{slide_0060.jpg}{Slide at 38:57}

So, now I'm going to give you the same thing as a homomorphism, which is a homomorphism modulo p, that is a homomorphism modulo p, that is a homomorphism from the invertible numbers modulo p to the group negative one, one, that's a multiplicative group. And for that you have this commutative diagram, where basically the decryption algorithm is, or the downward arrow, is applying the Legendre symbol, and going rightward is just multiplication. So, I'll just give you some letters as an example. So, if you start off with c and d in the upper left, those are numbers modulo p, you could immediately apply the Legendre symbols and go downward, and then you multiply them together, and because of the multiplicative property of the Legendre symbol, you get this equality that, what you end up with there is the Legendre symbol of c times d. So, but you could also take the Northeast passage, and just simply multiply c and d together to get c times d, and then at that point, apply the Legendre symbol. So, that's the diagram. So, how do we create a secure encryption scheme out of this? I mean, certainly p can't be known, because anyone could compute the Legendre symbol. \timestamp{39:17}


\section{Classical Schemes \& Quantum Attacks (40:00 -- 50:00)}

\keyslide{slide_0061.jpg}{Slide at 39:36}

So a possible answer, and the answer of Goldwasser-Macaulay, is that you hide P by using a number N, which is a product of P and Q, which are unknown. And then you have this quadratic residuosity assumption that given N and not P and Q, that it's hard to distinguish whether a number X modulo N that happens to have Jacobi symbol 1 is a square or a non-square. Okay, so this is kind of the mathematical way to creating an encryption scheme. We had a homomorphism, and we wanted to sort of harden it, and we did that by kind of masking it with this other prime Q. So now I want to describe the actual encryption scheme using sort of the encryption technique of Rothblum instead of the actual Goldwasser-Macaulay scheme. \timestamp{40:00}

\keyslide{slide_0062.jpg}{Slide at 40:16}

\keyslide{slide_0063.jpg}{Slide at 40:56}

So what you could do is you could have the public key be a bunch of Xs and Ys that encrypt either 1 or minus 1. So you just put lots of those in the public key. And then the encryption process is basically exactly what happens in Rothblum. You pick a random subset of 1 through K, and you take a product of the associated ciphertext from the public key to encrypt the value that you want. And decryption is just an application of the LeGenre symbol, modulo P. \timestamp{40:49}

\keyslide{slide_0064.jpg}{Slide at 41:35}

So, of course, this is not the way that Goldwasser-Macaulay does it. My point is sort of that they could have done it this way. They could have used this very generic technique of Rothblum to use homomorphism to create the magic security. The fact that they didn't need to do this was just, you know, kind of an artifact of some simplifications that could be made in this particular scheme. But I just kind of want to emphasize that the Rothblum technique is a very generic way to help create a semantically secure encryption scheme out of homomorphism. \timestamp{41:38}

\keyslide{slide_0065.jpg}{Slide at 42:15}

All right, so the shortcomings of Goldwasser-Macaulay, as well as other schemes like Elgamal and Pallier, are that they are not fully homomorphic, first of all. And they're not post-quantum. Quantum kills these schemes. And, in fact, quantum kills all homomorphic encryption schemes that use abelian groups, which is kind of sad. And the attack is actually kind of simple. \timestamp{42:27}

\keyslide{slide_0066.jpg}{Slide at 42:55}

There's this algorithm by Watris that for abelian groups, and actually for solvable groups, which include abelian groups, his algorithm is able to compute the order of the group. So, if you're using a group homomorphic encryption scheme, the encryptions of zero or whatever the group identity is, zero or one, encryptions of zero form a subgroup of the ciphertext. And so, at that point, all you have to do is use Watris' attack to distinguish between the encryption of zero subgroup. And, you know, if you want to tell if something is an encryption of zero, you just see if it's in this smaller group, which you can compute the order of, or if it's in the larger group. All right. \timestamp{43:16}

\keyslide{slide_0067.jpg}{Slide at 43:34}

\keyslide{slide_0068.jpg}{Slide at 44:14}

So, now I want to, you know, we saw that these schemes are not fully homomorphic. So, to get fully, full homomorphism, I want to move to ring homomorphisms, which allow two operations, not just one. So, both plus and times. Excuse me. So, examples of ring homomorphisms are modular reduction, just reducing an integer modular sum n, or a polynomial valuation, just, you know, so the evaluation of a product is the product of the evaluations, for example. \timestamp{44:05}

\keyslide{slide_0069.jpg}{Slide at 44:53}

And the idea is that the decryption algorithm would be some secret homomorphism that would commute with both plus and times. So, right away, there's a problem with this approach, namely, that quantum strikes again. So, if you look at the encryptions of zero, if decryption really is this ring of homomorphism, then the encryptions of zero form an ideal inside the ring of ciphertext. So, it's the kernel of homomorphism, and, in particular, it's an additive subgroup of the set of ciphertexts, of the group of ciphertexts, and Watchers' order-finding attack on group supplies again. So, there's no hope of this being post-quantum. \timestamp{44:55}

\keyslide{slide_0070.jpg}{Slide at 45:33}

Worse is that, basically, there's almost always a trivial linear algebra attack, where you just view the set of ciphertexts as a vector space, and the encryptions of zero will form a subspace, and it's very easy to use linear algebra to distinguish between the two. Worse is that, maybe there's some way to patch the approach, well, that's what we're going to try to do. So, back in 1978, Rivest, Edelman, and Jertuzos had a proposal, which was basically a ring-homomorphic approach, using a composite modulus, where the secret was a factor p. And, the encryption was just reducing modular p. So, then the question is how to hide p, and they do it here by using n, but that's not really enough for security in this setting, not for semantic security, because encryptions of zero are all multiples of p. \timestamp{45:44}

\keyslide{slide_0071.jpg}{Slide at 46:13}

Okay, so, in their scheme, a message m is basically encrypted as some number that's equal to m, modulo p. So, encryptions of zero are actual multiples of p, and just by taking the GCD of them, you can recover the secret key. So, they were well aware of the chosen plain text attacks back then, but they were still nice enough to propose these schemes, and we should all be thankful for it, because it's actually very similar to schemes that, scheme that I'll discuss next that's used today. Okay, so, this scheme, one idea for patching the scheme is that, maybe it gives a one-way encryption scheme, and maybe it could work if the plain text have enough entropy. So, that's basically what we'll do. \timestamp{46:33}

\keyslide{slide_0072.jpg}{Slide at 46:52}

\keyslide{slide_0073.jpg}{Slide at 47:32}

So, we'll try to hide the secret p by sort of adding noise or entropy to the messages to defeat GCD. And, what you end up with is this assumption, the approximate GCD assumption. That, if you sample lots of numbers x, that are each very close to multiples of p, but not exact multiples of p, then those integers, the assumption is, are indistinguishable from random integers of the same size. That's the approximate GCD assumption. And, we can use that to create a homomorphic encryption scheme. \timestamp{47:22}

\keyslide{slide_0074.jpg}{Slide at 48:11}

So, here's the scheme. The secret key is a large integer p, and then we'll create the public key, a la Rothblum, starting with a random approximate GCD instance, which, again, is just these integers x that are all really close to multiples of p. I'm saying that their residue modulo p is an even number, just for the sake of simplicity. And then, what we do to create the public key is we have these y's and z's, where y's are just the x, initial x's, they have even residues modulo p. And, the z's are, you add one to x's so that they have odd residues modulo p. \timestamp{48:11}

\keyslide{slide_0075.jpg}{Slide at 48:51}

And then, to encrypt, you just use rothbloom, you pick a random subset of these ciphertexts, such that the number of z's that you include, when you add them up, is odd if you want to encrypt one, and is even if you want to encrypt zero. And then, to decrypt, you first of all reduce modulo p, which should give you this value, right? And then, you take that value and you reduce that modulo two. And what you should get is the cardinality of s, which should be the message that was encrypted. So, for the homomorphism, you just add and multiply these integer ciphertexts together. \timestamp{49:00}

\keyslide{slide_0076.jpg}{Slide at 49:31}

And what happens is that, that, that implicitly adds and multiplies the, uh, the mod p residues together, which implicitly adds and multiplies the mod two values, uh, when you reduce modulo two, uh, \timestamp{49:50}


\section{Noise Management \& LWE (50:00 -- 60:00)}

\keyslide{slide_0077.jpg}{Slide at 50:10}

So security is based on Rothblum's procedure and the approximate GCG assumption. And so two problems that we have here is that the noise is growing. And this equation here that C mod P equals this two R stuff, that equation might not hold if the R's are bigger than P. And they wrap modular P. So at that point, you would have a decryption error. A second problem is that as you multiply these integers together, you get super big integers. \timestamp{50:00}

And so the size of the ciphertext in terms of the number of bits of the ciphertext is exploding. OK, so I'll skip this slide. I want to give you another example of a ring homomorphic approach using polynomial evaluation. So here we'll let R be the ring of integers modulo Q. And the homomorphism delta is just evaluating a polynomial at some secret point S. \timestamp{50:12}

OK, so all right, we're starting with this homomorphism. We want to create an encryption scheme, a homomorphic encryption scheme. How do we make it secure? OK, well, we have to hide S, obviously. But if this is but this doesn't really work because the encryptions of zero are all polynomials with a common root S. \timestamp{50:24}

\keyslide{slide_0078.jpg}{Slide at 50:50}

I mean, it's often said that computing Grubner bases is difficult. But in this case, if you don't want the number of monomials in your ciphertext to grow, if you have like a manageable number of monomials in your ciphertext, then you can just produce a vector space over these monomials. And so you'll have kind of a space, a vector space of ciphertext polynomials. \timestamp{50:36}

And the encryptions of zero will form a subspace within that space. And the linear algebra can easily distinguish encryptions of zero from encryptions of other things. So fellows in Koblitz back in 1993 had this, basically this scheme as the basic encryption scheme. They didn't really suggest homomorphic operations because of this problem that I mentioned, that as you add and multiply ciphertext, \timestamp{50:48}

you get a huge explosion of monomials. And so it becomes inefficient. So we want to repair this scheme. How do we do it? We're going to add noise and entropy to defeat linear algebra. \timestamp{51:00}

And here we sort of naturally arrive at the learning with errors assumption, which in terms of polynomials says this. It says that linear polynomials, fi, fi x, that evaluate to something very small, \timestamp{51:12}

\keyslide{slide_0079.jpg}{Slide at 51:30}

modulo q, at the secret s, are indistinguishable from random linear polynomials that might evaluate to something large of this. Okay. \timestamp{51:24}

So using this assumption, we produce the Pukorsky-Vikot-Tanathan homomorphic encryption scheme, where the secret key is just a random point. The public key, you can do a la Rothblum, \timestamp{51:36}

basically producing g's and h's as polynomials that evaluate to small numbers, which are even for encryptions of zero and odd for encryptions of one. And then we just use homomorphic addition \timestamp{51:48}

\keyslide{slide_0080.jpg}{Slide at 52:09}

of a random subset of those ciphertexts in the public key to produce a new ciphertext, whose parity, the parity of its noise is equal to the message that you want to encrypt. \timestamp{52:00}

And decryption is, well, you just evaluate the ciphertext polynomial, first of all, at s, modulo q. \timestamp{52:12}

And what should happen is that it's equal to this expression over the e's. And if that's true, then you can reduce that, modulo 2, \timestamp{52:24}

\keyslide{slide_0081.jpg}{Slide at 52:49}

to recover the cardinality of s, which should be n. All right. So you have the homomorphism is just addition and multiplication \timestamp{52:36}

of the ciphertext polynomials, which does the right thing to the messages inside them. Okay. So like the integer scheme, \timestamp{52:48}

we have these similar problems. We have, first of all, that the noise grows. And secondly, that the, \timestamp{53:00}

as you multiply ciphertext polynomials, the degree goes up. They're not linear polynomials anymore. So how do we fix that problem? Well, first of all, \timestamp{53:12}

\keyslide{slide_0082.jpg}{Slide at 53:28}

here's how we re-linearize. We take, so we have two linear ciphertext polynomials here. And what we want to do is multiply them together. \timestamp{53:24}

And when we do that, we get a quadratic polynomial. So hopefully, if you just evaluate that quadratic polynomial at s, \timestamp{53:36}

you get the right thing, which is e1 times e2 modulo q, which when you reduce that modulo 2, you get m1 times m2. But the ciphertext polynomial is quadratic. \timestamp{53:48}

\keyslide{slide_0083.jpg}{Slide at 54:08}

We want to eliminate the quadratic terms. How do we subtract them off? Well, what we do is we augment the public key with slightly quadratic polynomials. \timestamp{54:00}

Those are polynomials that are, that have, are mostly linear, except that they have one quadratic term. For example, \timestamp{54:12}

this polynomial has a quadratic term, two to the t times xj times xk. Okay. And you can use those slightly quadratic polynomials to subtract off \timestamp{54:24}

\keyslide{slide_0084.jpg}{Slide at 54:48}

all of the quadratic terms from the, from your c polynomial. So you just subtract all off all of the quadratic terms. \timestamp{54:36}

And what you end up with is a ciphertext that is linear again, and that encrypts the right thing because I didn't really mention it, \timestamp{54:48}

but the, but the, but the errors here are all even so that when you subtract off \timestamp{55:00}

\keyslide{slide_0085.jpg}{Slide at 55:27}

these, these polynomials, you preserve the parity of the error modulo 2 so that it encrypts the right thing. \timestamp{55:12}

Okay. And the second problem is bootstrapping. So you have a ciphertext polynomial \timestamp{55:24}

whose noise is getting kind of big. And if you perform any more operations on it, it's going to wrap modular Q and you're going to get a description error. \timestamp{55:36}

How do you reduce the error back to a manageable size? So one idea is to just try to reduce the noise by subtracting \timestamp{55:48}

\keyslide{slide_0086.jpg}{Slide at 56:07}

well-chosen encryptions at zero. You don't change what's encrypted, but this doesn't really work because you don't know how to subtract off the right noise. \timestamp{56:00}

You don't know the noise inside the ciphertext. So it doesn't, we haven't really made that approach work. \timestamp{56:12}

A second idea is that, well, you have the decryption algorithm and decryption and the seeker key are supposed to remove noise. \timestamp{56:24}

\keyslide{slide_0087.jpg}{Slide at 56:46}

That's the whole point. Is there some way that we can use them to reduce the noise level inside the encryption scheme? \timestamp{56:36}

And so here's how you, how that works. So here's your commutative diagram that we want for our encryption scheme. \timestamp{56:48}

And in fact, this commutative diagram works for some family of functions F. It's not all arithmetic circuits yet. \timestamp{57:00}

\keyslide{slide_0088.jpg}{Slide at 57:26}

It's kind of somewhat homomorphic. It works for, you know, polynomials up to a certain degree and then the noise becomes too big \timestamp{57:12}

and it doesn't work anymore. So it works for some family of functions F though. And our goal is to produce some fresh ciphertext \timestamp{57:24}

that encrypts the same message. So we start off with some ciphertext C that encrypts a message M and we want to kind of refresh \timestamp{57:36}

the ciphertext so that it has less noise. Okay. And we want to use our commutative diagram. \timestamp{57:48}

\keyslide{slide_0089.jpg}{Slide at 58:06}

This is where the commutative diagram really shines to do this. So we want to complete our commutative diagram. So how are we going to do it? \timestamp{58:00}

So what are we going to put in the lower left corner? Well, I claim that since the message M is hidden and this bootstrapping process \timestamp{58:11}

of refreshing a ciphertext is public and that the person that's doing the refreshing doesn't know anything about M \timestamp{58:24}

\keyslide{slide_0090.jpg}{Slide at 58:45}

except that C is the ciphertext that encrypts M. Basically, the only thing you can put here to complete this commutative diagram \timestamp{58:36}

is C because C is the only thing that I really know about M. Okay. So once I put C there, \timestamp{58:48}

that implies that this function F must be the decryption function because that's the only way I know to get from C to M. So, \timestamp{59:00}

\keyslide{slide_0091.jpg}{Slide at 59:25}

and once I've committed to C here, up here, we have that the ciphertext in the upper left. \timestamp{59:12}

These must be encryptions of, this is the decryption arrow, right? So these must be encryptions of the bits \timestamp{59:24}

of the ciphertext. And once, and once we, and then once we apply this evaluate function \timestamp{59:36}

using the function F, which we've determined to be the decryption function to this and these ciphertexts, \timestamp{59:48}


\section{Future of FHE (60:00 -- 65:22)}

\keyslide{slide_0092.jpg}{Slide at 60:05}

What we get here is a ciphertext C*, which is the evaluate function applied to the decryption function inside applied to the ciphertext X\_i. And so I'm not going to compute it exactly this way, I'll change it slightly and say that the decryption function, we don't know the secret key, right? So we can't evaluate exactly this function, instead we can consider it as the decryption \timestamp{60:00}

\keyslide{slide_0093.jpg}{Slide at 60:44}

function, which we know how to evaluate with an encrypted secret key and encrypted ciphertext bits. And that way we get basically the same thing, which is a ciphertext C*, which encrypts M if it's the case that this decryption function is in the family of functions F that the scheme is capable of evaluating. \timestamp{60:50}

\keyslide{slide_0094.jpg}{Slide at 61:24}

So this commutative diagram, everything is sort of predetermined. Once you decide that you want to refresh the ciphertext by exploiting this commutative diagram, which is essentially the essence of the homomorphic encryption scheme, you want to use the diagram and complete it, essentially everything is determined. And it seems like bootstrapping is sort of, in some sense, the only process that really works \timestamp{61:41}

\keyslide{slide_0095.jpg}{Slide at 62:03}

\keyslide{slide_0096.jpg}{Slide at 62:43}

to refresh a ciphertext. Okay. So I think I'm running over time at this point, but let's see. So what I want to say is that we finally achieved this commutative diagram. And if you consider the ciphertext space as an algebraic structure, it's very weird. \timestamp{62:32}

\keyslide{slide_0097.jpg}{Slide at 63:23}

It doesn't really preserve any of the original ring structure of the plaintext space. It has binary operations on it. And the binary operations are actually commutative, but they're not associative or anything else. So it has a structure of what's called a magma. So it's kind of maximally unstructured, aside from the commutativity, among algebraic structures. \timestamp{63:23}

\keyslide{slide_0098.jpg}{Slide at 64:02}

And one kind of wonders whether that's inherent in fully homomorphic encryption schemes that the ciphertext space and the operations that are performed through the evaluation algorithm have to be very unstructured in this sense. Okay. So I didn't really leave time for discussing new homomorphisms to explore. I'll try to zip through these slides in two minutes. \timestamp{64:14}

\keyslide{slide_0099.jpg}{Slide at 64:42}

As I said, one sort of wonders whether the ciphertext inherently needs to be incredibly unstructured to get a fully homomorphic encryption scheme. So first of all, we saw a real big problem when the ciphertext space was a solvable group. There were quantum attacks on that and non-solvable groups. Non-solvable groups are kind of interesting because you might apply You might get a homomorphism kind of using Barrington's theorem where Barrington had a theorem where he showed how to use non-abelian groups \timestamp{65:05}

to kind of evaluate circuits and you can do something kind of similar to that. So non-solvable groups are interesting in that sense, but there are all sorts of problems with using them. The quantum subgroup attack might still apply. We might just take an abelian subgroup of these non-solvable groups and still be able to distinguish between the big group and the subgroup. And there are all sorts of other problems. \timestamp{65:55}

Another problem is that usually these groups are linear. A linear group basically means a matrix group. It's isomorphic to a matrix group. And representation theory is basically the consideration of matrix groups. matrix groups and it's very powerful. \timestamp{66:46}

Break groups, for example, where, I mean, first of all, they were broken in lots of little ways, but they were broken definitively because, uh, uh, they were found to be linear, uh, using representation theory. And, uh, yeah, so there are all sorts of problems. Um, I still hold out hope for, uh, maybe a new FHE scheme based from multivariate crypto, where basically your, your ciphertext space is some sort of magma, where basically maybe you're doing something like, \timestamp{67:37}

uh, you're performing operations by doing some sort of weird inversions, you know, and, and, and, um, you're applying F and F inverse, which is, you know, something that's commonly done in multivariate systems. Um, and so that you would kind of hop around in the ciphertext space kind of pseudorandomly, like kind of like the, the last space schemes do. Um, but I haven't really fully explored that. \timestamp{68:28}

So it would, but it would be cool to have a, uh, HFE, uh, FHE scheme, uh, just, just for the acronym at least. Um, so in conclusion, uh, there's been great progress, uh, hopefully future breakthroughs. And, uh, thank you. \timestamp{69:19}


\end{document}
